\section{Тема 1. Линейные алгоритмы}

\subsection{Постановка задачи (вариант~4)}

Заданы моменты начала и конца некоторого промежутка времени
в часах, минутах и секундах (в пределах одних суток).
Найти продолжительность этого промежутка в тех же единицах
измерения.

\subsection{Математическое обоснование}

Каждый момент времени переводится в общее число секунд от начала
суток:
\[
  t = h \cdot 3600 + m \cdot 60 + s.
\]
Продолжительность промежутка вычисляется как
\[
  \Delta t = t_{\text{конец}} - t_{\text{начало}}.
\]
Если $\Delta t < 0$ (конец «раньше» начала из-за перехода через
полночь), к результату добавляется число секунд в сутках:
\[
  \Delta t' = \Delta t + 86400.
\]
Обратное преобразование в часы, минуты, секунды:
\[
  h = \left\lfloor \frac{\Delta t'}{3600} \right\rfloor,\quad
  m = \left\lfloor \frac{\Delta t' \bmod 3600}{60} \right\rfloor,\quad
  s = \Delta t' \bmod 60.
\]

\subsection{Блок-схема алгоритма}

\begin{center}
\begin{tikzpicture}[
  node distance=10mm, start chain=going below,
  nodes={draw=black, top color=white, bottom color=lightgray!50,
         minimum width=5.2cm, align=center, font=\small},
  arrow/.style={->, >=Stealth, thick},
]
  \node[draw, rounded corners, on chain] {Start};
  \node[shape=trapezium, trapezium left angle=70,
        trapezium right angle=110, on chain, draw, join=by arrow]
       {$h_1,m_1,s_1,\ h_2,m_2,s_2$};
  \node[shape=rectangle, on chain, draw, join=by arrow]
       {$t_1 := h_1{\cdot}3600+m_1{\cdot}60+s_1$};
  \node[shape=rectangle, on chain, draw, join=by arrow]
       {$t_2 := h_2{\cdot}3600+m_2{\cdot}60+s_2$};
  \node[shape=rectangle, on chain, draw, join=by arrow]
       (dt) {$\Delta t := t_2-t_1$};
  \node[shape=diamond, aspect=2, on chain, draw, join=by arrow]
       (cond) {$\Delta t < 0$\,?};
  \node[shape=rectangle, draw, right=2cm of cond]
       (fix) {$\Delta t := \Delta t+86400$};
  \node[shape=rectangle, on chain, draw, below=14mm of cond]
       (calc) {$\Delta h,\Delta m,\Delta s$ из $\Delta t$};
  \node[shape=trapezium, trapezium left angle=70,
        trapezium right angle=110, on chain, draw, join=by arrow]
       {$\Delta h$ ч $\Delta m$ мин $\Delta s$ с};
  \node[draw, rounded corners, on chain, join=by arrow] {End};
  \draw[arrow] (cond.east) -- node[above,font=\scriptsize]{$+$} (fix.west);
  \draw[arrow] (fix.south) |- (calc.east);
  \draw[arrow] (cond.south) -- (calc.north);
\end{tikzpicture}
\end{center}

\subsection{Текст программы}

\lstinputlisting[style=Cstyle]{../src/t01.c}

\subsection{Тестовые данные}

\begin{center}
\begin{tabular}{ll|l}
\toprule
Начало & Конец & Продолжительность \\
\midrule
10:30:15 & 14:45:50 & 04 ч 15 мин 35 с \\
23:50:00 & 00:10:00 & 00 ч 20 мин 00 с (переход через полночь) \\
08:00:00 & 08:00:00 & 00 ч 00 мин 00 с \\
\bottomrule
\end{tabular}
\end{center}

\textbf{Результат выполнения программы} (первый тест):
\begin{lstlisting}[language={},basicstyle=\ttfamily\small,frame=single]
Начало (чч мм сс): 10 30 15
Конец  (чч мм сс): 14 45 50
Продолжительность: 04 ч 15 мин 35 с
\end{lstlisting}
